\section{A complete proof of the rational-grid rank laws}
\label{rank:section}

The derivative-tower draft states an all-denominator rank law but supports it
only by finite exact computations.  The antiderivative draft states the
corresponding finite-range result carefully and leaves its extension to all
denominators as a conjecture.  Both questions admit the same elementary
character-theoretic solution.  The result concerns the rank of an explicitly
specified \emph{formal presentation}; it does not establish linear
independence of numerical zeta or Stieltjes values.

The relevant classical context is the universal ordinary distribution.
Kubert proved, in particular, that its level-$q$ module is free of rank
$\varphi(q)$ and becomes the regular unit-group representation over
$\mathbb Q$ \cite{Kubert1979}.  The proof below treats arbitrary rational
prime weights, states the endpoint convention explicitly, and obtains the
two rank formulas needed by the drafts.  We make no claim that the
underlying universal-distribution principle is new.

\subsection{The presentation, including the endpoint}

Fix $q\geq1$ and set
\[
 X_q=\tfrac1q\mathbb Z/\mathbb Z,
 \qquad V_q=\bigoplus_{x\in X_q}\mathbb Q e_x.
\]
When symbols are evaluated by a special function, use the representative
$\langle x\rangle_+\in(0,1]$.  In particular $e_0$ represents its value at
$1$, not its singular limit at $0$.  Choose a rational weight $w_p$ for each
prime $p\mid q$ and put
\[
 w(d)=\prod_{p\mid d}w_p^{v_p(d)}\qquad(d\mid q),
 \qquad w(1)=1.
\]
For every $d\mid q$ and $x\in X_{q/d}$ define a distribution row
\begin{equation}\label{rank:row}
 D_{d,x}=\sum_{\substack{y\in X_q\\dy=x}}e_y-w(d)e_x.
\end{equation}
There are exactly $d$ preimages in this sum.  Let $R_q(w)$ be the span of
these rows and let $U_q(w)=V_q/R_q(w)$.  The numerical application uses
$w_p=p^s$, with $s=k+1$ for positive parameter derivatives and $s=-k$ for
negative zeta jets.  Thus all the coefficient matrices are rational.

There are two distinct conventions worth recording.  The presentation
$U_q(w)$ retains the endpoint as a free symbol.  If the endpoint is placed
in a prescribed right-hand-side inventory, the homogeneous coefficient
space is
\begin{equation}\label{rank:endpoint}
 U_q^0(w)=V_q/(R_q(w)+\mathbb Q e_0).
\end{equation}
This is equivalent to moving every occurrence of the endpoint to the
right-hand side and using only the $q-1$ columns $e_{r/q}$,
$1\leq r<q$.

\begin{lemma}[Prime rows suffice]\label{rank:prime}
The rows $D_{p,x}$, for $p\mid q$ prime and $x\in X_{q/p}$, generate
$R_q(w)$.
\end{lemma}
\begin{proof}
For $ab\mid q$ and $x\in X_{q/(ab)}$, partition the $ab$ preimages by
their images under multiplication by $a$.  This gives the exact identity
\[
 D_{ab,x}=\sum_{\substack{y\in X_{q/a}\\by=x}}D_{a,y}
             +w(a)D_{b,x}.
\]
The factorization $w(ab)=w(a)w(b)$ holds even when $a$ and $b$ are not
coprime.  Repeated factorization proves the assertion, including all
prime-power rows.
\end{proof}

\subsection{The weighted character theorem}

Write $G_q=(\mathbb Z/q\mathbb Z)^\times$, acting by $u e_x=e_{ux}$.
For $q=1$ this is the trivial group.  If $\chi$ is a character of $G_q$,
write $f=f_\chi\mid q$ for its conductor and $\chi^*$ for its primitive
character.  The principal character has conductor $1$ and
$\chi^*(p)=1$.  For every $m$ with $f\mid m\mid q$, put
\begin{equation}\label{rank:E}
 E_{m,\chi}=
 \sum_{a\in(\mathbb Z/m\mathbb Z)^\times}
       \chi^*(a)e_{a/m}.
\end{equation}
At $m=1$ this is $e_0$.  These vectors are used after extending scalars to
$\mathbb C$; they need not individually have rational coefficients.

\begin{theorem}[Weighted distribution module]\label{rank:main}
For every $q\geq1$ and every choice of rational weights $w_p$, the
complexification of $U_q(w)$ has the basis
\[
 \bigl\{[E_{f_\chi,\chi}]:\chi\in\widehat G_q\bigr\}.
\]
The $\chi$-labelled basis vector has $G_q$-character $\chi^{-1}$.
Consequently
\begin{equation}\label{rank:dimensions}
 \dim_{\mathbb Q} U_q(w)=\varphi(q),
 \qquad
 \dim_{\mathbb Q} U_q^0(w)=\varphi(q)-1.
\end{equation}
In particular the first dimension is independent of all prime weights.
\end{theorem}

\begin{proof}
Decompose $V_q\otimes\mathbb C$ by the exact order $m\mid q$ of its
symbols.  The summand of order $m$ is the regular representation of
$G_m$, inflated to $G_q$.  The reduction map $G_q\to G_m$ is surjective
by the Chinese remainder theorem.  Hence a character $\chi$ of $G_q$
occurs in that summand if and only if $f_\chi\mid m$, and its
$\chi^{-1}$-eigenline is spanned by $E_{m,\chi}$.  This identifies every
character block before imposing relations.

By Lemma~\ref{rank:prime}, it is enough to sum prime distribution rows
against $\chi^*(a)$ over symbols of each exact order $m$.  Suppose
$pm\mid q$ and $f\mid m$.  If $p\mid m$, all $p$ lifts of a unit
$a/m$ have exact denominator $pm$.  Each unit modulo $pm$ occurs once,
with its correct character coefficient.  The resulting row is
\begin{equation}\label{rank:edge-ramified}
 E_{pm,\chi}-w_p E_{m,\chi}.
\end{equation}
If $p\nmid m$, exactly one lift has denominator $m$: writing that lift
as $b/m$, one has $pb\equiv a\pmod m$.  Its coefficient is
$\chi^*(a)=\chi^*(p)\chi^*(b)$.  The other lifts run through the units
modulo $pm$.  In this case the row is
\begin{equation}\label{rank:edge-unramified}
 E_{pm,\chi}-(w_p-\chi^*(p))E_{m,\chi}.
\end{equation}
The source space for rows of exact order $m$ has no $\chi$-component
when $f\nmid m$.  Since the distribution map is $G_q$-equivariant,
there are no further rows in the $\chi$-block.  This also explains why
a newly appearing primitive character is not removed when its conductor
is first reached.

For $f\mid m\mid q$ define
\begin{equation}\label{rank:A}
 A_{m,\chi}=
 \prod_{p\mid f}w_p^{v_p(m)-v_p(f)}
 \prod_{\substack{p\mid m\\p\nmid f}}
       (w_p-\chi^*(p))w_p^{v_p(m)-1}.
\end{equation}
Empty products and zeroth powers have value $1$, so $A_{f,\chi}=1$,
also if some weight is zero.  Multiplying the edge coefficients along
any divisor chain from $f$ to $m$ gives exactly $A_{m,\chi}$:
different prime steps commute.  Thus every class in this block satisfies
\begin{equation}\label{rank:transport}
 [E_{m,\chi}]=A_{m,\chi}[E_{f,\chi}],
\end{equation}
and the block has dimension at most one.

For the opposite inequality, define on the block the linear functional
$E_{m,\chi}\mapsto A_{m,\chi}$.  Formula~\eqref{rank:A} shows directly
that it vanishes on every edge row
\eqref{rank:edge-ramified}--\eqref{rank:edge-unramified}; its value on
$E_{f,\chi}$ is $1$.  The block therefore has dimension exactly one.
Summing over all $\varphi(q)$ characters proves the first formula in
\eqref{rank:dimensions}.  Rank does not change under the scalar extension
$\mathbb Q\subset\mathbb C$.  Finally $e_0=E_{1,\chi_0}$ is the nonzero
generator of the principal block, so adjoining $e_0$ to the relation
space removes exactly that block and proves the second formula.
\end{proof}

\begin{remark}[Resonance affects a choice of basis, not the rank]
\label{rank:resonance}
For $w_p=p^s$ with real $s\ne0$, every factor in
\eqref{rank:A} is nonzero, since $|\chi^*(p)|=1$ when $p\nmid f$.
For nonrational real weights the same proof is read over $\mathbb R$
instead of $\mathbb Q$, followed by extension to $\mathbb C$.
The classes of the $\varphi(q)$ unit symbols $e_{a/q}$ then form a basis
of $U_q(w)$: their Fourier coordinates are the nonzero multiples
$E_{q,\chi}=A_{q,\chi}E_{f,\chi}$.  At resonant weights this particular
basis can fail.  For example, if $w_p=1$ and $q=p$ is prime,
$E_{p,\chi_0}=0$ whereas $e_0$ still spans the principal block.
The conductor basis in Theorem~\ref{rank:main} remains valid, so there
is no exceptional rank.  No division by $w_p-\chi^*(p)$ occurs in the
proof.
\end{remark}

\subsection{The full finite-jet module is free}

The same argument controls all jet layers at once.  This gives more than
adding the dimensions of successive layers: it proves a free presentation
before the nilpotent jet parameter is specialized.

\begin{theorem}[Freeness over a coefficient algebra]\label{rank:algebra}
Let $K$ be a field of characteristic zero containing the values of all
characters of $G_q$, and let $A$ be any commutative $K$-algebra with
identity.  Define $V_q(A)$ and its distribution rows as in
\eqref{rank:row}, now with arbitrary weights $w_p\in A$ and with
$A$-linear spans.  Then
\[
 V_q(A)/R_q(w;A)\simeq A^{\varphi(q)},\qquad
 V_q(A)/(R_q(w;A)+Ae_0)\simeq A^{\varphi(q)-1}.
\]
The conductor vectors $E_{f_\chi,\chi}$ give the first basis; omitting
the principal vector gives the second.  These statements remain valid
when $A$ has zero divisors or nilpotents and when edge coefficients vanish.
\end{theorem}
\begin{proof}
The Fourier changes of basis over each $G_m$ and the character
projections divide only by $|G_m|$, which is invertible in $K$ and
therefore in $A$.  The decomposition into character blocks thus works
over $A$ exactly as it does over $\mathbb C$.
Within a block, the edge equations
\eqref{rank:edge-ramified}--\eqref{rank:edge-unramified} express every
generator as $A_{m,\chi}$ times the conductor generator.  They give a
surjection $A\to U_\chi$, $1\mapsto[E_{f_\chi,\chi}]$.
The map $E_{m,\chi}\mapsto A_{m,\chi}$ descends to $U_\chi\to A$
and is its inverse, since it sends the conductor generator to $1$.
Hence each block is free of rank one, without any use of division by
a weight or an Euler factor.  Removing the principal block gives the
endpoint assertion.
\end{proof}

\begin{corollary}[Finite Hurwitz and Stieltjes jets]\label{rank:full-jets}
Fix an integer $s\ne1$ and $N\geq0$, and let $K\subseteq\mathbb C$
contain all character values modulo $q$ and all $\log p$, $p\mid q$.
For $0\leq j\leq N$ use the formal symbols
$Z_j(x)=\zeta^{(j)}(s,\langle x\rangle_+)/j!$; the equality here
specifies their intended evaluation, not an arithmetic-independence
assumption.  Impose every coefficient of Hurwitz distribution through
order $N$, and prescribe all $N+1$ endpoint symbols $Z_j(0)$.
The resulting formal $K$-vector space has dimension
\begin{equation}\label{rank:full-jet-dimension}
 (N+1)(\varphi(q)-1).
\end{equation}
Before prescribing the endpoint its dimension is $(N+1)\varphi(q)$.
For fixed $k\geq1$, the same result holds for the complete finite
Stieltjes-derivative array
$\{\gamma_j^{(k)}(r/q):0\leq j\leq N,\ 1\leq r<q\}$,
with all its coefficient distribution rows and all endpoint layers
prescribed.
\end{corollary}
\begin{proof}
Use the truncated polynomial algebra $A=K[\eta]/(\eta^{N+1})$ and
the weights
\[
 w_p=p^s\sum_{\ell=0}^N\frac{(\log p)^\ell}{\ell!}\eta^\ell.
\]
They give $w(d)=d^s\exp(\eta\log d)$ in this algebra.  The coefficient
row at order $j$ is
\begin{equation}\label{rank:coefficient-jet-row}
 \sum_{dy=x}Z_j(y)
 -d^s\sum_{\ell=0}^j
       \frac{(\log d)^\ell}{\ell!}Z_{j-\ell}(x).
\end{equation}
Identify $Z_j(x)$ with $\eta^{N-j}e_x$ in the underlying $K$-vector
space of $V_q(A)$.  This reversed grading is necessary: multiplication
by $\eta$ lowers jet order.  Under this identification,
\eqref{rank:coefficient-jet-row} is exactly $\eta^{N-j}D_{d,x}$,
and prescribing all endpoint layers is exactly quotienting by $Ae_0$.
Theorem~\ref{rank:algebra} proves freeness over $A$; multiplying its
rank by $\dim_K A=N+1$ proves \eqref{rank:full-jet-dimension}.

Finally the coefficient identity
\[
 \sum_{j=0}^N\frac{(-1)^j}{j!}\gamma_j^{(k)}(a)\eta^j
 =(-1)^k(1+\eta)_k
       \sum_{j=0}^N\frac{\zeta^{(j)}(k+1,a)}{j!}\eta^j
       \pmod{\eta^{N+1}}
\]
is an invertible triangular change of jet coordinates: its scalar
factor has nonzero constant coefficient $(-1)^k k!$.  It preserves
the coefficient distribution presentation, with weights
$d^{k+1}\exp(\eta\log d)$, and proves the Stieltjes assertion.
\end{proof}

In particular, passing from an individual highest layer to the entire
finite array introduces no additional formal rank obstruction.  The
freeness result concerns these displayed relations over the specified
coefficient field; it still gives no numerical-independence theorem.

\subsection{Reflection and the negative-jet conjecture}

For $\epsilon\in\{1,-1\}$, let $S_q^\epsilon$ be generated by
$e_x+\epsilon e_{-x}$ for $x\in X_q$.  On the $\chi$-block, reflection
$x\mapsto-x$ acts by $\chi(-1)$, so these rows remove exactly the
characters satisfying $\chi(-1)=\epsilon$.

\begin{corollary}[Distribution and reflection rank]\label{rank:reflection}
For every $q\geq3$, every rational choice of weights, and
$\epsilon\in\{1,-1\}$,
\begin{equation}\label{rank:reflection-dimension}
 \dim_{\mathbb Q}
 \frac{V_q}{R_q(w)+S_q^\epsilon+\mathbb Q e_0}
 =\frac{\varphi(q)}2-\frac{1-\epsilon}{2}.
\end{equation}
For $q=1,2$ this dimension is $0$ for either sign.
\end{corollary}
\begin{proof}
The surviving characters have parity $-\epsilon$.  For $q\geq3$,
$-1$ is a nonidentity element of order two in $G_q$, so half of its
characters have each parity.  The principal character survives the
reflection quotient exactly when $\epsilon=-1$, and is then removed
by $e_0$.  For $q=1,2$, $G_q$ has only the principal character, which
the endpoint relation removes in either case.
\end{proof}

Here is the precise application to the drafts.  Put
$J_k(x)=\zeta'(-k,\langle x\rangle_+)$ for $k\geq1$.  Differentiating
Hurwitz multiplication gives
\begin{equation}\label{rank:jet-distribution}
 \sum_{j=0}^{d-1}\zeta'\left(-k,\frac{a+j}{d}\right)
 -d^{-k}\zeta'(-k,a)
 =d^{-k}\log d\,\zeta(-k,a).
\end{equation}
The reflection row used in the antiderivative draft has left-hand side
$J_k(x)+(-1)^kJ_k(-x)$.  More explicitly, for $0<x<1$ and
$C_\nu(x)+iS_\nu(x)=\operatorname{Li}_{\nu}(e^{2\pi i x})$, its known
right-hand side is
\begin{equation}\label{rank:jet-reflection}
 J_k(x)+(-1)^kJ_k(-x)
 =\begin{cases}
  \displaystyle\frac{(-1)^{k/2}k!}{(2\pi)^k}C_{k+1}(x),&k\text{ even},\\[6pt]
  \displaystyle\frac{(-1)^{(k-1)/2}k!}{(2\pi)^k}S_{k+1}(x),&k\text{ odd}.
 \end{cases}
\end{equation}
Indeed, differentiate the standard Hurwitz Fourier formula at $s=-k$;
the corresponding cosine or sine factor has a simple zero, so its
derivative is the sole surviving factor.  The endpoint value
$J_k(0)=\zeta'(-k)$ is placed in the inventory separately.

Take as the inventory the endpoint $\zeta'(-k)$, all Bernoulli-value
terms $d^{-k}\log d\,\zeta(-k,a)$ in
\eqref{rank:jet-distribution}, and all right-hand sides of
\eqref{rank:jet-reflection}.  Then
Corollary~\ref{rank:reflection}, with $s=-k$ and $\epsilon=(-1)^k$,
proves for \emph{all} $q\geq3$ and $k\geq1$ that the formal residual
dimension is
\begin{equation}\label{rank:negative-result}
 r_{q,k}=\begin{cases}
  \varphi(q)/2-1,&k\text{ odd},\\
  \varphi(q)/2,&k\text{ even}.
 \end{cases}
\end{equation}
This proves the unrestricted version of the antiderivative draft's
Proposition \texttt{prop:jetrank}, formerly checked only for
$3\leq q\leq30$ and $1\leq k\leq4$.  Equivalently, the remaining
complex character lines are indexed by primitive characters of
conductors $f\mid q$, $f>1$, with $\chi(-1)=(-1)^{k+1}$.

For the positive parameter tower, let the inventory contain
$\gamma_n^{(k)}(1)$ and every lower-$s$-derivative-layer term occurring
in the differentiated multiplication identity.  At a fixed highest
layer its homogeneous coefficient is $d^{k+1}$, independent of $n$.
In particular the $n=1$ row is
\[
 \sum_{j=0}^{d-1}\gamma_1^{(k)}\left(\frac{a+j}{d}\right)
 -d^{k+1}\gamma_1^{(k)}(a)
 =d^{k+1}\log d\,\psi^{(k)}(a).
\]
Theorem~\ref{rank:main} therefore proves the formal residual dimension
$\varphi(q)-1$ for every $q\geq1$ and $k\geq1$.  For $n>1$ this is
an associated-layer statement: lower derivatives are prescribed data,
not additional unknown columns.

\paragraph{What the theorem does and does not certify.}
The background inventory is part of the statement.  Several of its
terms, including rational-argument polygamma and Clausen values, are
not known to be elementary.  One may define an evaluation map from the
formal quotient to the span of actual constants modulo the inventory;
the theorem determines the dimension of its source.  Its image can
have smaller dimension, and proving injectivity would be a separate
arithmetic-independence problem.  Consequently ``one surviving formal
character line'' cannot be replaced by ``one numerically independent
new constant.''  The same caution applies to unsuccessful PSLQ searches.

\paragraph{Exact reproducibility.}
The accompanying \texttt{code/verify\_ranks.py} forms the actual
$q$-column rational matrices, including prime and all-divisor rows,
endpoint rows, and both reflection signs.  It checks their ranks by
exact fraction elimination and records its parameters and results in
\texttt{data/rank\_checks.json}.  It also checks the resonant weight
$s=0$ and a non-power weight system.  These computations are regression
checks on the presentation and its boundary conventions; the proof for
unbounded $q$ and $k$ is Theorem~\ref{rank:main}.
The companion \texttt{code/verify\_jet\_ranks.py} checks the expanded
finite-jet coefficient matrices with rational exponential-weight
parameters and with nilpotent weights; its output is
\texttt{data/jet\_rank\_checks.json}.
